% page 622 du fac-simile (image p-621 du scan)
% bloc 152.4 x 246.3 mm ; marge gauche 8.0 mm
\pgc{-12.700mm}{0.000mm}{
\l{\cel{4}614}
\l{}
\l{}
\l{versiono Maniero ye qua persono komprenas lekte, na-}
\l{\cel{3}racas, raportas, explikas ulo. DEFIS.}
\l{}
\l{versoro. (matem.)Lineo di qua la longeso e la situeso}
\l{\cel{3}karakterizas nombro imaginara. -\cel{2}EF.}
\l{}
\l{vertebro. (anat.) Singla ek la 24 osti ye qui konsistas}
\l{\cel{3}la spino. - DEFIS.}
\l{}
\l{vertico.\cel{2}(\sou{anat.}) La suprajo di la kapo. - EFIS.}
\l{}
\l{vertijar. (\sou{netrans,}) (\sou{patol.})\cel{2}Subisar ta sturdo kurta-}
\l{\cel{3}dura dum qua semblas ke la objekti jiras o rotacas}
\l{\cel{3}avan la okuli. - EFIS.}
\l{}
\l{\sou{vertikala}. (\sou{geom.)}\cel{3}Perpendikla ye la plano horizontala.}
\l{\cel{3}- DEFIRS.}
\l{}
\l{\sou{vertuo}. Praktiko kustumata di lo etiko-bona. - EFIS.}
\l{}
\l{\sou{veruko}. I. Mikra surkreskajo pelala, surfacala. - II.}
\l{\cel{3}Protuberanco sur la kortico di arboro, sur folio.}
\l{\cel{3}- FIS.}
\l{}
\l{\sou{vervo}.\cel{2}Inspireso ardoroza, nur kelka-dura. - DEF.}
\l{}
\l{\sou{vespo}. (\sou{zool.})\cel{2}Insekto, genero ek la tribuo \textquotedbl{}himenop-}
\l{\cel{3}teri\textquotedbl{}, di qua la femino (qua havas dardo longa, quale}
\l{\cel{3}la abelo) konstruktas alveoli. - L.\cel{1}\sou{vespa}. - DEIS.}
\l{}
\l{\sou{vespero}. La parto finala di la jorno-periodo. - EFL.}
\l{}
\l{\sou{vespertilio}. (\sou{zool.)}\cel{2}Mikra mamifero, karnavida, noktala,}
\l{\cel{3}di qua la osti di la membri avana unionesas a la kor-}
\l{\cel{3}po di la membri dopa per membrano quaze ala. - L.\cel{1}\sou{ves}-}
\l{\cel{3}\sou{pertilio}. - L.}
\l{}
\l{\sou{vespri}. (\sou{liturgio katol.}) Recitodi la oficio deala quan}
\l{\cel{3}onu agis olim cirkum vespero, e nun posdimeze. - DEFIS.}
\l{}
\l{\sou{vesto}. To quo uzesas por kovrar la korpo homala, skope di}
\l{\cel{3}shirmar lu de la domaji da vetero.mala, o cesigar lua}
\l{\cel{3}nudeso. - EFIS.}
\l{}
\l{\sou{vestalo}. (\sou{historio di la Roma antiqua}). Sacerdotino di}
\l{\cel{3}Vesta, qua devis permanar virga vivo-dure, e diqua la}
\l{\cel{3}rolo esis entratenar la fairo sakra.- (\sou{metaf.}) Muliero}
\l{\cel{3}tre chasta.- DEFIRS.}
\l{}
\l{\sou{vestio}. Vesto sen baski. - F.}
\l{}
\l{vestibulo. I. Chambro qua, situita ye la enireyo di edifico,}
\l{\cel{3}di domo, igas acesar olca. - II. Placo avan la pordo}
\l{\cel{3}precipua di kirko o katedralo. - DEFIS.}
}
